% !TeX program = xelatex
% ==========================================================================
% Пропоузал семестрового проєкту - курс «Штучний інтелект», УКУ.
% Компіляція: XeLaTeX + BibTeX (Overleaf робить це автоматично).
% Локально:   latexmk (див. README.md шаблону).
% ==========================================================================
\documentclass[ukrainian]{ucuproposal}

% --------------------------------------------------------------------------
% Метадані титульної сторінки
% --------------------------------------------------------------------------
\projecttitle{Тематичне моделювання telegram пропаганди: порівняння середовищ і динаміки тем}
\projecttagline{Пайплайн, що перетворює кластери постів на структуровані твердження й порівнює про- та антикремлівські канали до і після 24.02.2022}

\teammember{Онищук Артем}
\teammember{Полянський Іван}
\teammember{Лещук Роман}

\mentor{ Косован Олександр, Корець Христина} % якщо є ментор - впишіть ПІБ

\track{research}

\repo{} % додати посилання на репозиторій до camera-ready

\proposaldate{2026-09-20}

\begin{document}
\maketitlepage

\section{Проблема та мотивація}
Тематичні карти воєнного Telegram будували вже не раз. Проблема в тому,
що з кластера на тисячу постів і десятка ключових слів ще не видно, що саме канал стверджує. Ми хочемо
довести пайплайн далі, до твердження виду «актор, дія, обʼєкт, оцінка»,
яке вже можна порахувати, порівняти між середовищами й відстежити в часі.

На вході у нас пости прокремлівських, антикремлівських і білоруських
проурядових каналів: текст, дата, канал, перегляди й метадані форвардів.
На виході - теми з твердженнями та їхні частки по середовищах і тижнях, а
з них уже можна відповісти чим відрізняються ландшафти двох таборів, що
зʼявилося або змінило вагу після 24.02.2022, які теми спалахують навколо
подій, а які тримаються фоном, як швидко тема розходиться каналами і чи
збігаються найчастіші теми з найпереглядуванішими. Розмітки при цьому
немає, тексти короткі й шумні, а надійної автоматичної метрики якості тем
теж немає.

\section{Огляд джерел та наявних рішень}
LDA~\cite{blei2003lda} на коротких постах працює погано, тому для
соцмереж стандартом став BERTopic~\cite{grootendorst2022bertopic}, який
поєднує контекстні ембединги, UMAP~\cite{mcinnes2018umap},
HDBSCAN~\cite{campello2013hdbscan} і c-TF-IDF. На воєнному Telegram його
вже застосовували~\cite{tretiakov2024telegram}, тут автори у~\cite{bawa2025battlefield,bawa2024dissent} аналізували теми
антикремлівського сегмента за фазами вторгнення, а в білоруському
датасеті~\cite{rudnik2026autocrats} уже лежать 137 тем від Top2Vec,
названих за допомогою LLM. Тому карта тем вже зроблена давно, тому
ми беремо її як вхід і надбудовуємо рівень тверджень і
стенсу~\cite{hardalov2022stance} та динаміку тем~\cite{kleinberg2002bursts}.
Окремо про оцінювання: Hoyle та ін.~\cite{hoyle2021coherence} показали,
що автоматична когерентність погано валідована для нейронних
топік-моделей і часто оголошує переможця там, де люди різниці не бачать.
Ми будемо використовувати LLM як анотатора, але з перевіркою на
людях~\cite{gilardi2023chatgpt}.

\section{Дані}
Основний корпус - це \textit{Pro-/Anti-Kremlin Telegram
Dataset}~\cite{bawa2025data} під CC~BY~4.0: 404 прокремлівські канали
(4\,109\,645 постів) і 114 антикремлівських (1\,117\,768) за рік до і рік
після вторгнення.

Білоруський корпус~\cite{rudnik2026autocrats}, теж CC~BY~4.0, менший: 20
каналів і 23\,570 постів за лютий-березень 2025 року. З основним корпусом
він не перетинається в часі жодним днем, тому прямого порівняння
«Білорусь проти Росії» робити не будемо - воно показало б різницю
новинних циклів, а не середовищ. Використаємо його як out-of-domain тест:
перенесемо навчену модель і подивимося, що саме переноситься, звіряючись
із опублікованими темами Top2Vec.

Мова в усіх трьох середовищах переважно російська, зокрема і в
білоруських каналах. Виняток - пʼять зі 114 антикремлівських каналів
пишуть українською (УНІАН, СБУ, «Київ Оперативний»). Це, насправді трохи проблема, бо модель може ділити табори за мовою, а не за змістом, тому проганяємо language ID і звітуємо метрики з
цією підвибіркою і без неї.

Архіви важать 4,8~ГБ, тому розпаковуємо не все: вибірку на 300 тис.
постів відбираємо потоково, рівномірно по каналах і місяцях, без
дублікатів і форвардів. Як запас тримаємо резервний
датасет~\cite{kyrychenko2025data}. Сирі дані не комітимо, ліцензії
фіксуємо в \texttt{data/README.md}.

\section{Підхід і методи}
Пайплайн має 3 рівні на реалізаціях BERTopic,
sentence-transformers, gensim.
\textit{Кластеризація.} Основна модель - BERTopic, одна
на обʼєднаному корпусі, щоб теми різних середовищ і періодів лежали в
спільному просторі. Головне питання тут - чи варто взагалі перекладати:
порівнюємо багатомовні ембединги на оригіналі з
англомовною моделлю поверх машинного перекладу тієї ж підвибірки.

\textit{Твердження.} Далі кластер перетворюємо на твердження за схемою
«актор, дія, обʼєкт, оцінка» плюс стенс щодо України, Заходу і влади: LLM
отримує ключові слова й близько 20 постів кластера і повертає JSON,
невалідні відповіді перезапитуємо.

\textit{Динаміка й порівняння.} Будуємо тижневі ряди часток тем, шукаємо
сплески за алгоритмом Клейнберга~\cite{kleinberg2002bursts} і ділимо теми
на подієві та фонові. Дивимось скільки
минає часу від появи теми в каналі-джерелі до того, як її підхоплять
інші. Відстань між середовищами рахуємо дивергенцією Дженсена-Шеннона.
Рахуємо на безкоштовній GPU у Colab/Kaggle, seed фіксуємо.

\section{Оцінювання результатів}
Автоматичну частину рахуємо для всіх конфігурацій проти LDA: NPMI,
$C_v$~\cite{roder2015coherence}, topic diversity і стабільність (ARI між
запусками з різними seed). Звітуємо їх, але вирішальними не
вважаємо~\cite{hoyle2021coherence}.

Також буде людська оцінка. Word intrusion і topic intrusion на 40
темах, троє анотаторів незалежно.
Для тверджень збираємо золотий набір зі 150 кластерів і міряємо точність
по кожному полю схеми, F1 по стенсу і частку тверджень, які анотатор не
може підтвердити постами кластера. Детектор сплесків звіряємо з датами
реальних подій.

На виході буде репозиторій із пайплайном, таблиця метрик, золотий набір і
звіт.

\section{План і ризики}
\textbf{До мідтерму}: потокова вибірка, EDA, language ID, LDA, перша
модель BERTopic, експеримент із перекладом і попередні відповіді про
різницю таборів та зсув після 24.02.2022. \textbf{До фіналу}: схема
тверджень і золотий набір, динаміка й поширення, перенос на білоруський
корпус, людська оцінка і звіт.

Найбільший ризик - ручна розмітка, бо не знаємо скільки часу займе. Якщо не вкладаємося, ріжемо золотий набір до 80 кластерів і не
чіпаємо автоматичну частину. Якщо LLM не тримає схему на шумних
кластерах, спрощуємо її до пари «актор, оцінка». Якщо не вистачить
обчислень, зменшуємо вибірку й беремо легшу ембединг-модель.

\section{Contribution statement}
\begin{contributions}
  \contribution{Онищук Артем}{вибірка й препроцесинг, EDA, language ID, LDA baseline}
  \contribution{Полянський Іван}{BERTopic-пайплайн, експеримент із перекладом, метрики якості та стабільності}
  \contribution{Лещук Роман}{схема тверджень і LLM-екстракція, динаміка й поширення тем, білоруський корпус}
\end{contributions}
Розмітку золотого набору, людську оцінку тем і звіт робимо разом. Усі
учасники згодні з таким розподілом.

\ucuprintbibliography

\end{document}
